\documentclass[12pt, twoside]{book}
\usepackage[letterpaper, includehead, bottom=2.5cm, left=3cm, right=3cm, top=2.5cm, headsep=1.25cm]{geometry}
\usepackage[spanish]{babel}
\usepackage{graphicx}
\graphicspath{{./images/}}
\usepackage[table]{xcolor}
\usepackage{fancyhdr}
\usepackage{lmodern}
\usepackage{enumerate, multicol}
\usepackage{enumitem}
\usepackage{tocloft}
\usepackage{amssymb}

\usepackage{emptypage}
\usepackage{tikz}
\usepackage[most]{tcolorbox}
\usepackage{imakeidx}
\usepackage{eso-pic}
\usepackage{titlesec}
\usepackage{colortbl}
\usepackage{csquotes}
\usepackage{booktabs}
\usepackage{subcaption}
\usepackage{hyperref}
%% configs and definitions


%colors
\definecolor{ustgreen}{HTML}{00685a}
\definecolor{ustmidgreen}{HTML}{009b87}
\definecolor{ustlightgreen}{HTML}{00b49d}
\definecolor{ustorange}{HTML}{ff6d3c}
\definecolor{ustdarkgreen}{HTML}{004d43}
\definecolor{graytitle}{HTML}{69827f}
\definecolor{grayheader}{HTML}{5c706d}
\definecolor{darkgraytitle}{HTML}{151a19}
\definecolor{lightbackground}{HTML}{bae7e2}
\definecolor{ustblue}{HTML}{006db6}
\definecolor{ustlightblue}{HTML}{69c0d3}
\definecolor{ustdarkblue}{HTML}{003d6b}
\definecolor{ustskyblue}{HTML}{4fa8c2}
\definecolor{usttealblue}{HTML}{006f8e}
\definecolor{ustred}{HTML}{d94c4c}
\definecolor{ustdarkred}{HTML}{9e1b32}
\definecolor{ustlightred}{HTML}{ff8074}
\definecolor{ustbrickred}{HTML}{9b3d3d}
\definecolor{ustscarlet}{HTML}{f44e42}
\definecolor{tablemid}{RGB}{232,251,247}
\definecolor{tabletop}{RGB}{192, 204, 202}

% basic definitions
\newcommand{\set}[1]{\left\lbrace #1 \right\rbrace}

%format of the title for chapters
\titleformat{\chapter}[display]
{\bfseries\large}
{\LARGE\textcolor{graytitle}{Unidad \thechapter}\\* \titlerule}
{2pt}
{\huge \centering\color{darkgraytitle}\titlerule[2pt]\vspace{0.85cm}}
[\vspace{0.85cm}{\titlerule[2pt]}]
\titlespacing*{\chapter}{0pt}{-40pt}{40pt}

% fancyhdr
\newcommand{\numberedchapter}{
    \cleardoublepage
    \fancyhead[RE]{{\bfseries \leftmark}}
    \fancyhead[LO]{{\bfseries \rightmark}}
}
\newcommand{\unnumberedchapter}[1]{
    \cleardoublepage
    \addcontentsline{t oc}{chapter}{#1}
    \fancyhead[RE]{{\bfseries #1}}
    \fancyhead[LO]{{}}
}
\pagestyle{fancy}
\fancyhf{}
\fancyhead[LE,RO]{\includegraphics[width=2.5cm]{logo_horizontal.png}}
\fancyhead[RE, LO]{\color{grayheader} Departamento de Ingeniería}
\fancyfoot[LO]{\small\color{grayheader} Dr. José Alejandro Aburto}
\fancyfoot[LE,RO]{\color{grayheader} \thepage}


% hyperref setup
\hypersetup{
    colorlinks=true,
    linkcolor=ustdarkblue,
    filecolor=black!30!magenta,
    urlcolor=black!40cyan,
    pdfauthor={Dr. José Alejandro Aburto Araneda},
}

%index
\makeindex[columns=2]

% thm & tcolorboxes numbering
\newtheorem{theorem}{Teorema}[chapter]
\newtcolorbox[auto counter, number within=chapter]{definition}[2][]{colback=ustblue!2, colframe=ustblue, fonttitle=\bfseries, title=Definición~\thetcbcounter \hfill \textcolor{ustblue!10}{ #2 #1}}
\newtcolorbox[auto counter, number within=chapter]{example}[1][]{colback=ustgreen!5, colframe=ustgreen, fonttitle=\bfseries, title=Ejemplo~\thetcbcounter, breakable, label=#1}
\newtcolorbox[auto counter, number within=chapter]{examples}[1][]{colback=ustgreen!5, colframe=ustgreen, fonttitle=\bfseries, #1, title=Ejemplos~\thetcbcounter, breakable}
\newtcolorbox{note}{enhanced, colback=ustred!7, frame style={left color=ustdarkred, right color=ustred}, fonttitle=\bfseries, sharp corners, title=¡Cuidado!, watermark tikz={\draw[line width=2mm] circle (1cm)
node{\fontfamily{ptm}\fontseries{b}\fontsize{20mm}{20mm}\selectfont !};}}
\newtcolorbox{simpleframe}{colframe=ustbrickred, float, enhanced}
\newtcolorbox[auto counter, number within=chapter]{thm}[2][]{
    enhanced,
    breakable,
    colback=usttealblue!7,
    frame style={
        left color=usttealblue,
        right color=usttealblue!70
    },
    fonttitle=\bfseries,
    title=Teorema~\thetcbcounter \hfill \textcolor{usttealblue!13}{#1} #2
}
\newtcolorbox[auto counter, number within=chapter]{proposition}[2][]{
    enhanced,
    breakable,
    colback=usttealblue!5,
    frame style={
        left color=usttealblue!80,
        right color=usttealblue!60
    },
    fonttitle=\bfseries,
    title=Proposición~\thetcbcounter \hfill \textcolor{usttealblue!13}{#1} #2
}

\newtcolorbox[auto counter, number within=chapter]{observation}[1][]{
    enhanced,
    breakable,
    colback=ustlightgreen!7,
    frame style={
        left color=ustlightgreen,
        right color=ustlightgreen!90
    },
    fonttitle=\bfseries,
    title=Observación~\thetcbcounter \hfill \textcolor{ustlightgreen!13}{#1}
}

\begin{document}
% titlepage
\begin{titlepage}
    \newgeometry{left=7.5cm}
    \noindent
    \begin{flushright}
        \color{white}{\fontsize{32}{54} \selectfont \textsf{Matemáticas para Ciencias e Ingeniería I}}
    \end{flushright}

    \makebox[0pt][l]{\rule{1.3\textwidth}{1pt}}
    \par
    \pagecolor{ustmidgreen}

    \noindent {\Large \textcolor{ustdarkgreen}{\textsf{Universidad Santo Tomás}}\ \textcolor{ustlightgreen!10}{\textsf{Facultad de Ingeniería}}}
    \vfill
    \noindent
    \textcolor{ustlightgreen!15}
    {\large \textsf{Departamento de Ingeniería}}
    \vskip\baselineskip
    \noindent
    \textcolor{white}{\textsf{2025}}
\end{titlepage}

\restoregeometry
\nopagecolor

\pagenumbering{roman}
%\thispagestyle{empty}
\tableofcontents

\chapter*{Presentación}
\addcontentsline{toc}{chapter}{Presentación}
Este apunte busca ser una guía esencial y concisa para el curso de primer año \emph{Matemáticas para Ciencias e Ingeniería I} que comparten diversas carreras en la Universidad Santo Tomás.


Incluye el contenido en pocas palabras, incluyendo ejemplos y ejercicios propuestos. En versiones posteriores se pretende agregar:
\begin{enumerate}[label=\roman*)]
    \item Ejercicios resueltos, con mayor complejidad que los ejemplos básicos del contenido
    \item Agregar un solucionario para los ejercicios
    \item Añadir una sección o subsección en cada capitulo con ideas o problemas interesantes para profundizar
    \item Añadir una bibliografía
\end{enumerate}
\vfill
\noindent
Matemáticas para Ciencias e Ingeniería © 2025 by Dr. José Alejandro Aburto is licensed under CC BY-NC-SA 4.0. To view a copy of this license, visit https://creativecommons.org/licenses/by-nc-sa/4.0

\pagenumbering{arabic}
\chapter{Lógica y Polinomios}
\section{Lógica proposicional}
Comencemos con una definición fundamental. Queremos formalizar el concepto de \emph{proposición}, las cuales trabajaremos usualmente como incógnitas y operaremos con operaciones binarias (que requieren dos) semejantes a la suma o al producto usuales.
\begin{definition}{Proposición}
    Una proposición es una expresión que puede poseer solo uno de dos valores: \textbf{verdadero} ó \textbf{falso}.
\end{definition}
Podemos comenzar pensando en qué no es una proposición, y hay muchísimos ejemplo: todas las preguntas no son proposiciones. Primero tenemos proposiciones en el lenguaje cotidiano, y también expresiones matemáticas.
\begin{example}
    \begin{enumerate}
        \item Todos los alumnos de la UST deben sacar \(5.5\) o más para eximirse.
        \item El número \(5\) es par.
        \item La ecuación \(x^2+x+1\) no tiene solución en los números reales.
        \item Todos los gatos son grises.
    \end{enumerate}
\end{example}


Tendremos tres operaciones:
\begin{itemize}[label=\(\bullet\)]
    \item La conjunción, \enquote{y}, que denotaremos por  \(\wedge\)
    \item La disyunción, \enquote{ó}, que denotaremos por \(\vee\)
    \item La negación, \enquote{no}, que denotaremos por \(\neg\)
\end{itemize}
\begin{note}
Hay 3 notaciones para la negación, por ejemplo:
    \[\neg p=\sim p =\overline{p}\]
\end{note}
\section{Tablas de verdad}
Como las proposiciones solo pueden tener dos valores: \textbf{verdadero} ó \textbf{falso}, podemos describir los resultados de las operaciones al usarlas en una tabla. Este tipo de tabla les llamaremos \textbf{tablas de verdad}. Las siguientes tablas son por definición:
\begin{table}[h]
    \centering
    \begin{subtable}{0.3\linewidth}\centering
            \begin{tabular}{ccc}
                \toprule
                \(p\)   &   \(q\)   &   \(p\wedge q\) \\
                \midrule
                \(V\)   &   \(V\)   &   \(V\) \\
                \rowcolor{tablemid}
                \(V\)   &   \(F\)   &   \(F\) \\
                \(F\)   &   \(V\)   &   \(F\) \\
                \rowcolor{tablemid}
                \(F\)   &   \(F\)   &   \(F\) \\
            \end{tabular}
            \caption{Conjunción}
    \end{subtable}
    \hfill
         \begin{subtable}{0.3\linewidth}\centering
            \begin{tabular}{ccc}
                \toprule
                \(p\)   &   \(q\)   &   \(p\vee q\) \\
                \midrule
                \(V\)   &   \(V\)   &   \(V\) \\
                \rowcolor{tablemid}
                \(V\)   &   \(F\)   &   \(V\) \\
                \(F\)   &   \(V\)   &   \(V\) \\
                \rowcolor{tablemid}
                \(F\)   &   \(F\)   &   \(F\) \\
            \end{tabular}
            \caption{Disyunción}
    \end{subtable}
    \hfill
    \begin{subtable}{0.3\linewidth}\centering
            \begin{tabular}{cc}
                \toprule
                \(p\)   &   \(\overline{p}\) \\
                \midrule 
                \(V\)   &   \(F\)   \\
                \rowcolor{tablemid}
                \(F\)   &   \(V\)
            \end{tabular}
            \caption{Negación}
\end{subtable}
\end{table}

\begin{note}
    Las disyunción no es exclusiva, es decir, si decimos \enquote{\(p\) ó \(q\)} pueden ocurrir ambas a la vez. Hay que prestar atención a esto al principio ya que en el lenguaje cotidiano suele ser exclusivo, o sea, solemos elegir solo uno de los dos.
\end{note}


Podemos utilizar las tablas de verdad para determinar los valores posibles de una proposición compuesta, es decir, una proposición construida con estos operadores usando otras proposiciones. Por ejemplo:
\begin{example}
    Consideremos la proposición compuesta: \((p\vee q)\wedge (\overline{p}\vee\overline{q})\). Sería difícil desarrollarlo en solo un paso, por lo que agregaremos columnas a la tabla de manera que nos ayude a calcular el resultado, poniendo partes más simples que forman la proposición compuesta que nos interesa:
    \begin{center}
        \begin{tabular}{cccc}
            \toprule
            \(p\)   &   \(q\)   &    \(p\vee q\)    &  \(\overline{p}\vee\overline{q}\) \\
            \midrule
            \(V\)   &   \(V\)   &   \(V\)           & \(F\) \\
            \rowcolor{tablemid}
            \(V\)   &   \(F\)   &   \(V\)           & \(V\)\\
            \(F\)   &   \(V\)   &   \(V\)           &\(V\)\\
            \rowcolor{tablemid}
            \(F\)   &   \(F\)   &   \(F\)           & \(V\)
        \end{tabular}
    \end{center}
Ahora que ya hemos calculado estas proposiciones más simples, podemos calcular los valores posibles de la proposición compuesta inicial:
    \begin{center}
        \begin{tabular}{ccccc}
            \toprule
            \(p\)   &   \(q\)   &    \(p\vee q\)    &  \(\overline{p}\vee\overline{q}\) & \((p\vee q)\wedge (\overline{p}\vee\overline{q})\)\\
            \midrule
            \(V\)   &   \(V\)   &   \(V\)           & \(F\) & \(F\)\\
            \rowcolor{tablemid}
            \(V\)   &   \(F\)   &   \(V\)           & \(V\) &\(V\)\\
            \(F\)   &   \(V\)   &   \(V\)           &\(V\)  &\(V\)\\
            \rowcolor{tablemid}
            \(F\)   &   \(F\)   &   \(F\)           & \(V\) & \(F\)
        \end{tabular}
    \end{center}
\end{example}
Además, tenemos otro operador que aparece usualmente en matemáticas, por ejemplo en cada resultado. Solemos tener una hipótesis que es aquello que suponemos como cierto y una consecuencia de esos resultados. Podemos pensar en el teorema de Pitágoras como ejemplo, partimos de un triángulo rectángulo en el plano, y como consecuencia obtenemos una ecuación que satisfacen sus lados.
\begin{definition}{Implica}
    Definimos el conectivo (u operador)	\textbf{implica} (o \textbf{entonces}) como se indica en la siguiente tabla:
    \begin{center}
        \begin{tabular}{ccc}
            \toprule
            \(p\)   &   \(q\)   &   \(p\Rightarrow q\) \\
            \midrule
            \(V\)   &   \(V\)   &   \(V\)\\
            \rowcolor{tablemid}
            \(V\)   &   \(F\)   &   \(F\)\\
            \(F\)   &   \(V\)   &   \(V\)\\
            \rowcolor{tablemid}
            \(F\)   &   \(F\)   &   \(V\)\\
        \end{tabular}
    \end{center}
\end{definition}

Notemos que si completamos la siguiente tabla de verdad:
\begin{table}[h]
    \centering
    \begin{tabular}{cccc}
          \toprule
          \(p\)   &   \(q\)   &   \(p\Rightarrow q\)  & \(\overline{p}\vee q\)\\
            \midrule
          \(V\)   &   \(V\)   &   \(V\)               & \(V\)\\
            \rowcolor{tablemid}
          \(V\)   &   \(F\)   &   \(F\)               & \(F\)\\
          \(F\)   &   \(V\)   &   \(V\)               & \(V\)\\
            \rowcolor{tablemid}
          \(F\)   &   \(F\)   &   \(V\)               & \(V\)\\

    \end{tabular}
    \caption{Primera equivalencia lógica}
\end{table}
    Las dos últimas columnas tienen los mismos valores de verdad pero no son la misma proposición. Por ejemplo, \enquote{dos es un número par} y \enquote{5 es un número primo} tienen el mismo valor de verdad pero no son la misma proposición. Por ello, son equivalentes de cierta forma pero no iguales. Esto motiva usar un símbolo diferente:
\begin{definition}{Equivalencia lógica}
    Cuando dos proposiciones (compuestas) \(R\) y \(T\) poseen los mismos valores (en el mismo orden) en la tabla de verdad, diremos que son \textbf{equivalentes} y lo denotaremos por:
    \[R\equiv T\]
\end{definition}
Esto es una relación como equivalencia al igual que la igualdad cuando trabajamos los números reales por dar un ejemplo.

\section{Álgebra Booleana}
Ahora que tenemos un sentido de \enquote{igualdad}, quisieramos poder desarrollar expresiones como en el álgebra habitual. A esta operatoria y sus propieades les llamamos \textbf{Álgebra Booleana}.

Tenemos varias propiedades que nos serán de utilidad y enlistamos ahora:
\begin{thm}{Propiedades del álgebra booleana}
    \begin{enumerate}[label=\roman*.]
        \item Conmutatividad
            \[p\wedge q\equiv q\wedge p \qquad p\vee q \equiv q\vee q\]
        \item 
        \item 
        \item 
    \end{enumerate}
\end{thm}

Esto nos permite trabajar con proposiciones compuestas de una manera más convienente muchas veces.
\begin{example}
    Se tienen 3 proposiciones: \(p\), \(q\), \(r\). Sabemos que \(p\vee q \equiv r \wedge q\), queremos desarrollar y simplificar la siguiente proposición compuesta: \(\left[(p\wedge q)\vee r\right]\vee \overline{r}\). Procedemos a desarrollar utilizando las propiedades anteriores:
    \begin{align*}
        \left[(p\vee q)\wedge r\right]\vee \overline{r} & \equiv \left[(r\wedge q)\wedge r\right]\vee \overline{r} \\
                & \equiv \left[(q\wedge r)\wedge r\right]\vee \overline{r} \\
                & \equiv \left[(q\wedge (r\wedge r)\right]\vee \overline{r} \\
                & \equiv \left[q\wedge r\right]\vee\overline{r} \\
                & \equiv (q\vee\overline{r})\wedge(r\vee\overline{r}) \\
                & \equiv (q\vee\overline{r})\wedge V \\
                & \equiv q\vee\overline{r}.
    \end{align*}
\end{example}
Aunque en general podríamos obtener un resultado solo en incógnitas, a veces podremos hallar el valor.
\begin{example}
    Tenemos que
    \begin{align*}
        (p\wedge q)\vee(p\Rightarrow\overline{q})
        & \equiv (p\wedge q)\vee(\overline{p}\vee \overline{q}) \\
        & \equiv (p\wedge q)\vee \overline{\left(p\wedge q\right)} \\
        & \equiv V
    \end{align*}
\end{example}

Ya sea con álgebra booleana o con las tablas de verdad, a veces ocurrirá que sin importar los casos una proposición compuesta sea verdadera, falsa o con un valor indeterminado. Esto tiene un nombre definido como sigue:
\begin{definition}{Tautología - Contingencia - Contradicción}
    Si una proposición compuesta da como resultado
        \begin{itemize}[label=\(\bullet\)]
                \item verdadero, le llamaremos \textbf{Tautología}
                \item falso, le llamaremos \textbf{Contradicción}
                \item indeterminado, le llamaremos \textbf{Contingencia}
        \end{itemize}
\end{definition}
\begin{note}
    El que una contingencia tenga un valor indeterminado significa que depende de los valores de \(p,q,r,\ldots\).
\end{note}

\section{Funciones Proposicionales}
\begin{definition}{Función proposicional}
    Una \textbf{función proposicional} es una función que toma como argumento un elemento de un universo definido y devuelve una proposición.
\end{definition}
\begin{example}
    Tenemos las siguientes funciones proposicionales:
    \begin{enumerate}[label=\roman*.]
        \item \(q(k)\equiv \text{\(k\) es par y mayor que \(5\)}\)
        \item \(F(x)\equiv x^2+x+1=0\)
        \item \(r(t)\equiv t \text{ es divisible por \(3\)}\)
    \end{enumerate}
\end{example}
En estos casos no hemos definido un universo, el cual es el dominio de la función. Sin embargo, si el contexto es claro podemos omitir esta información. 
\begin{observation}
    En general tendremos dos usos para las funciones proposicionales: Uno será con cuantificadores lógicos y el otro en la definción de conjuntos. En ambos casos estará escrito con claridad cuál es el universo en el que estaremos trabajando.
\end{observation}

\section{Cuantificadores Lógicos}
Para describir qué tanto una función proposicional nos proporciona proposiciones verdaderas, utilizaremos cuantificadores lógicos. 
\begin{definition}{Cuantificador Lógico}
    Un \textbf{cuantificador lógico} es un operador que se aplica a las funciones proposicionales para indicar qué tantas proposiciones resultan ser verdaderas. Tenemos tres cuantificadores:
    \begin{itemize}[label=\(\bullet\)]
        \item \textbf{Para todo / Para cada}: \(\forall\)
        \item \textbf{Existe alguno / Hay algún}: \(\exists\)
        \item \textbf{Existe un único}: \(\exists!\)
    \end{itemize}
\end{definition}
\begin{example}
    \begin{enumerate}
    \item \enquote{Todos los números enteros son pares ó impares} podemos escribirlo como:
        \[\forall n\in\mathbb{Z}: n\text{ es par} \vee n\text{ es impar}\]
    \item \enquote{Existe un número par} podemos escribirlo como:
        \[\exists n\in\mathbb{Z}: n\text{ es par}\]
    \item \enquote{Para cada número natural hay un número primo mayor que él} podemos escribirlo como:
        \[\forall n\in\mathbb{N},~ \exists p\in\mathbb{N}: p\text{ es primo} ~ \wedge ~ p>n\]
    \end{enumerate}
\end{example}
Para determinar el valor de verdad de cada una de las proposiciones compuestas en las que está involucrado un cuantificador lógico, tenemos que seguir los siguientes razonamientos:
\begin{observation}
    \begin{itemize}[label=\(\bullet\)]
        \item Para que un \textbf{para todo} sea verdad, necesitamos que todos los elementos del universo nos dé verdadero al pasarlos por la función proposicional, por ello, si tan solo uno de los elementos nos da como resultado una proposición falsa, el para todo será falso.

        \item Para que un \textbf{existe} sea verdadero bastará con que uno de los elementos del universo nos de verdadero, sin importar si hay varios que nos sirvan para esto, con uno basta. Por otro lado, si todos nos dan como resultado una proposición falsa, el existe será falso.

        \item Para que un \textbf{existe un único} sea verdadero, necesitaremos que haya un elemento del universo que nos de verdadero, y además que todos los demás sean falsos. Por lo tanto, el existe un único será falso si hay dos \emph{elementos diferentes} que nos dan una proposición verdadera, ó si todos los elementos del universo nos dan falso.
    \end{itemize}
\end{observation}
\begin{example}
    \begin{enumerate}
        \item \enquote{Todos los números enteros son pares} podemos escribirlo como:
            \[\forall n\in\mathbb{Z}: n\text{ es par}\]
        En este caso, esta proposición es \textbf{falsa}. Esto, debido a que existe un\footnote{si hay más no importa, solo nos interesa que haya al menos un elemento del universo que resulte en una proposición falsa para que el \textbf{para todo} sea falso.} número entero que no es par, particularmente \(n=3\).
    \item \enquote{Existe un único número real que satisface \(x^2-1=0\)} podemos escribirlo como:
            \[\exists! x\in\mathbb{R}: x^2-1=0\]
        En este caso, esta proposición es \textbf{falsa}. Esto, debido a que existen dos números reales que satisfacen la proposición, particularmente \({x=1}\) y \({x=-1}\).
    \end{enumerate}
\end{example}
Ahora quisieramos saber como resulta ser la negación de una de estas proposiciones compuestas en las que hay cuantificadores lógicos involucrados.


Aunque ya hemos hablado acerca de cómo saber una proposición de este tipo es verdadera o falsa, no hemos visto formalmente la negación de estas proposiciones:
\begin{proposition}{Negación de Cuantificadores Lógicos}
    Tenemos las siguientes maneras de negar una proposición compuesta con cuantificadores lógicos:
    \begin{itemize}[label=\(\bullet\)]
        \item La negación del \textbf{para todo}:
            \[\neg\left(\forall u\in U: P(u)\right)\equiv \exists u\in U: \neg P(u)\]
        \item La negación del \textbf{existe}:
            \[\neg\left(\exists u\in U: P(u)\right)\equiv \forall u\in U: \neg P(u)\]
        \item La negación del \textbf{existe un único}:
            \begin{align*}
            \neg\left(\exists!\ u\in U: P(u)\right) & \equiv \left(\forall u\in U: \neg P(u)\right) \\
                                                    & \vee \exists u_1,u_2\in U: P(u_1)\wedge P(u_2)\wedge u_1\neq u_2
            \end{align*}
    \end{itemize}
\end{proposition}

\begin{example}
    Consideremos la proposición siguiente:
    \[T:\text{\enquote{El cuadrado de cada número real es mayor que 1}}\]
    La podemos escribir formamente como:
    \[T: \forall x\in\mathbb{R}: x^2>1\]
    La negación de esta proposición es:
    \[\neg T: \exists x\in\mathbb{R}: x^2\leq 1\]
    Podemos observar que \(\neg T\) es una proposición verdadera ya que por ejemplo \(x=0\) cumple que \(x^2\leq 1\). Por lo tanto, \(T\) es falsa.
\end{example}


\section{Conjuntos}
Debido a que puede resultar ser complicado definir un conjunto de una manera formal lo entenderemos como una colección de elementos. Tenemos varias maneras de representarlos:
\begin{definition}{Conjunto}
    Un \textbf{conjunto} es una colección de elementos de un universo dado. Podemos representarlo mediante:
    \begin{itemize}[label=\(\bullet\)]
        \item Definición intensiva (o por compresión): Especificamos un universo, y una propiedad (una proposición) que los elementos deben cumplir.
        \item Definición extensiva: Especificamos todos los elementos explícitamente.
        \item Definición ostensiva: Especificamos ejemplos que representan al conjunto.
    \end{itemize}
\end{definition}
\begin{note}
    Solo la definición intensiva y la extensiva son formales. Esto debido a que la definición ostensiva puede llevar a errores dependiendo de la interpretación que le demos.
\end{note}
\begin{example}[example:definiciones_conjuntos]
    \begin{enumerate}
        \item Consideremos el conjunto:
            \[\{x\in\mathbb{R}: x^2-1<0\}\]
        Este conjunto está definido de manera intensiva ó por comprensión.
        \item Consideremos el conjunto:
            \[\{2,3,5,7,9,\pi\}\]
        Este conjunto está definido de manera extensiva.
        \item Consideremos el conjunto:
            \[\{2,3,5,7,\ldots\}\]
        Este conjunto está definido de manera ostensiva.
    \end{enumerate}
\end{example}
\begin{observation}
    En el ejemplo \ref{example:definiciones_conjuntos} notemos la manera en que está escrita el conjunto
    \({x\in\mathbb{R}: x^2-1<0}\). La definición se lee como:
    \begin{center}
        \enquote{Todos los reales \(x\) tales que \(x^2-1<0\)}
    \end{center}
    Aquí, los dos puntos \enquote{\(:\)} ó bien la barra vertical \enquote{\(|\)} se lee como \enquote{tal que}.
    

    Además, el símbolo \(\in\) se lee como \enquote{pertenece a}.
\end{observation}
\begin{note}
    La manera ostensiva es solamente una representación informal. Por ejemplo el conjunto \(\{3,5,7,\ldots\}\) podemos interpretarlo de al menos dos maneras. ¿Se trata de números primos mayores que 2 o de números impares mayores que 2?
\end{note}

\section{Operaciones entre Conjuntos y Diagrama de Venn}
Tenemos tres operaciones elementales, la \textbf{unión}, la \textbf{intersección} y el \textbf{complemento}. Cada uno se obtendrá de las tres operaciones en lógica.

Denotaremos el universo por \(U\) a menos que se indique lo contrario.
\begin{definition}{Operaciones entre Conjuntos}
    \begin{itemize}[label=\(\bullet\)]
        \item \textbf{Unión}: 
            \[A\cup B=\set{x\in U: x\in A\vee x\in B}\]
        \item \textbf{Intersección}: 
            \[A\cap B=\set{x\in U: x\in A\wedge x\in B}\]

        \item \textbf{Complemento}:
            \[U\setminus A=A^c=\set{x\in U: x\notin A}\]
        \item \textbf{Diferencia}\footnote{Notemos que el símbolo que utilizamos es como un menos pero inclinado. Así lo diferenciamos del menos usual}:
            \[A\setminus B=\set{x\in U: x\in A\wedge x\notin B}\]
    \end{itemize}
\end{definition}
Para facilitar la compresión veamos unos ejemplos concretos sencillos:
\begin{example}
    Consideremos los siguientes conjuntos:
    \[A=\set{-1,1,3,6,8},\qquad B=\set{1,5,8}\]
    \[U=\set{-1,0,1,2,3,4,5,6,7,8}\]
    Y así, tenemos que:
    \begin{enumerate}[label=\alph*)]
        \item La unión de \(A\) y \(B\) es contiene los elementos de ambos conjuntos:
            \[A\cup B=\set{-1,1,3,5,6,8}\]
        \item La intersección de \(A\) y \(B\) es el conjunto de los elementos que pertenecen a ambos conjuntos a la vez:
            \[A\cap B=\set{1,8}\]
        \item El complemento de \(A\) es el conjunto de todos los elementos que no pertenecen a \(A\):
            \[U\setminus A=\set{0,2,4,5,7}\]
        \item La diferencia de \(A\) y \(B\) es el conjunto de los elementos de \(A\) que no pertenecen a \(B\):
            \[A\setminus B=\set{-1,3,6}\]
    \end{enumerate}
\end{example}
\begin{note}
    La diferencia de \(A\) y \(B\) solo quita elementos a el conjunto \(A\). Si \(B\) tiene elementos que \(A\) no posee, son totalmente ignorados. Por ejemplo, si \({A=\set{1,2}}\) y \({B=\set{2,3,4,5}}\), obtendremos que \({A\setminus B=\set{1}}\).


Notar además que \(B\setminus A=\set{3,4,5}\).
\end{note}

%% diagramas
Podemos representar los conjuntos de una manera gráfica, especialmente cuando se trata de conjuntos abstractos ó si son de finitos elementos (no son infinitos). En el caso de los conjuntos abstractos, estos diagramas nos ayudarán a comprender mejor las operaciones entre conjuntos ya que podremos visualizarlas fácilmente.

\begin{definition}{Diagrama de Venn}
    Un \textbf{diagrama de Venn} es una representación gráfica de conjuntos que utiliza figuras geométricas cerradas para representar conjuntos. Particularmente,
\end{definition}

\section{Conjuntos numéricos}

\section{Polinomios}

\section{División de polinomios}
\section{Teorema del Resto}
\section{Ceros Reales y Racionales}

\section{Números complejos}

\chapter{Geometría Analítica}

\chapter{Trigonometría y Relación de Orden}

\chapter{Funciones y Límites}



\printindex

\begin{titlepage}
    \newgeometry{left=7.5cm}

    \makebox[0pt][l]{\rule{1.3\textwidth}{0pt}}
    \par
    \pagecolor{ustmidgreen}

    \noindent
\end{titlepage}

\nopagecolor
\end{document}
